\documentclass[conference]{IEEEtran}
\usepackage{amsmath,amssymb}
\usepackage{booktabs}
\usepackage{array}
\usepackage{url}
\usepackage{cite}

\title{Is There a Limit to Human Intelligence?\\
A Conditional Intelligence Frontier Benchmark for Measurement, Plasticity, and Neural Constraints}

\author{\IEEEauthorblockN{Research Design Report}
\IEEEauthorblockA{Four-Agent Scientific Workflow: Survey--Idea--ExperimentDesign--Author}}

\begin{document}
\maketitle

\begin{abstract}
The question of whether human intelligence has a limit is scientifically meaningful only after the target, task, age, environment, and measurement scale are specified. IQ is a norm-referenced score, not a physical unit, and a historical increase in test performance does not by itself establish an equal increase in a fixed biological quantity. This paper reframes the problem as a test of a conditional intelligence frontier: does performance in calibrated cognitive domains approach a reproducible asymptote, and can education, health, or structured training move that asymptote or only improve practiced tasks? We synthesize evidence on the Flynn effect, education, polygenic architecture, brain systems, and transfer. We propose a design-only benchmark combining high-range item-response calibration, alternate-form testing, a multi-age longitudinal cohort, a randomized adaptive-training module with an active control, and an EEG/MRI sub-study. A hierarchical latent-state model separates a general factor, domain factors, speed, practice, and measurement error. A limit is supported only when a near-zero trajectory slope is stable across difficult item forms, delayed follow-up, and independent sites. The forecast is conditional: domain-specific constraints and individual asymptotes are plausible, while a universal IQ maximum is not identified by current evidence. No participant, neural, causal, or intervention results are reported.
\end{abstract}

\begin{IEEEkeywords}
human intelligence, IQ, Flynn effect, cognitive plasticity, latent variable, psychometrics, neuroscience, learning curves
\end{IEEEkeywords}

\section{Introduction}
\label{sec:intro}
The claim that IQ scores rose by roughly thirty points during the twentieth century is a popular entry point to a deeper question: does human intelligence have an upper limit? The observation is not trivial. Intelligence-test performance has changed across generations, and these changes are large enough to alter the interpretation of norms and educational comparisons. However, the observation is also not a direct measurement of a universal biological quantity. An IQ score is calibrated against a reference population, depends on the selected test, and is sensitive to education, nutrition, health, language, motivation, and familiarity with formal testing.

The phrase human intelligence also combines distinct constructs. Fluid reasoning, acquired knowledge, working memory, processing speed, spatial reasoning, and learning a novel task are related but not identical. A person can show an age-related decline in speed while retaining knowledge, or improve a reasoning strategy without improving far-transfer performance. A common factor can summarize covariance among domains, but it should not erase domain-specific trajectories.

This paper therefore asks:
\begin{quote}
Under specified task domains, ages, environments, practice histories, and measurement conditions, does human cognitive performance approach a reproducible upper boundary, and how much of that boundary is shifted by education, health, training, and biological variation?
\end{quote}

The question has four requirements. First, an apparent plateau must be separated from a psychometric ceiling caused by insufficiently difficult items. Second, test-retest and training gains must be separated from generalization. Third, age and cohort effects must not be treated as interchangeable. Fourth, brain and genetic correlates must be tested as possible mediators or constraints rather than interpreted as immutable destiny.

The contribution is a design-only benchmark with a falsifiable decision structure. It combines evidence synthesis with a preregistration-ready human protocol. The protocol does not claim to discover an absolute maximum for every possible person, task, culture, or future technology. It can test whether a conditional frontier is stable under declared conditions and whether specific inputs move it.

\section{Survey: What Is Already Known?}
\label{sec:survey}

\subsection{The Flynn effect is a trend in measured performance}
The Flynn effect describes rising scores on many cognitive tests across generations. A formal meta-analysis covering 1909--2013 found worldwide gains but substantial heterogeneity by country, period, test, and cognitive domain \cite{pietschnig2015}. This is strong evidence that measured performance is environmentally sensitive. It is not evidence that a fixed latent factor, the upper tail of every ability, or the biological capacity for learning has increased by the same amount.

Several mechanisms can produce cohort gains. Education may increase reasoning strategies and knowledge. Improved nutrition and public health may reduce developmental constraints. Technological and occupational environments may exercise unfamiliar forms of abstraction. Changes in family size, schooling access, test familiarity, and migration can alter the sampled population. A trend can also reflect changes in the construct measured by an instrument. Therefore a cohort comparison requires measurement invariance: item difficulty, discrimination, and factor meaning should be sufficiently comparable across time and groups.

The opposite risk is to treat a slowdown or reversal in some national series as proof that a universal ceiling has been reached. A decline can reflect changing schooling, health, immigration, inequality, test composition, or selective participation. The scientific target is not the sign of an average trend alone, but the latent trajectory under a declared measurement model.

\subsection{Education and other modifiable inputs}
Education provides a candidate causal lever. A meta-analysis of quasi-experimental and other designs reported that additional education improves intelligence-test performance, with effects varying across outcomes and study designs \cite{ritchie2018}. This finding is compatible with a finite but movable frontier. A person may gain access to strategies and representations that move performance within a task family without making all cognitive operations unlimited.

Sleep, nutrition, disease burden, stress, language exposure, and cognitive opportunity can alter both performance and measurement reliability. Their effects can be immediate, developmental, or cumulative. They also interact: a person with the same latent capacity may express it differently under fatigue, sensory load, or a time limit. A design that samples only one testing session cannot distinguish a stable capacity from a state-dependent expression.

\subsection{Polygenic and distributed biological architecture}
Behavioral-genetic and molecular-genetic results describe intelligence differences as substantially heritable but highly polygenic. Plomin and von Stumm review a literature in which many variants of small effect contribute to individual differences and emphasize that heritability does not imply immutability \cite{plomin2018}. Large genome-wide association studies identify distributed associations rather than a single intelligence locus \cite{savage2018}. These results make a universal biological cutoff unlikely to be identifiable from one gene or one pathway.

Neuroscience similarly points to distributed systems. Reviews link intelligence differences to brain structure, information-processing characteristics, and network organization \cite{deary2010}. More recent work integrates genetic variation, imaging, and intelligence differences while stressing the complexity of the mapping \cite{deary2021}. Neural measures may help identify constraints, efficiency, reserve, or compensation, but a correlation with a score is not a proof of mechanism. A predictive brain model can also fail under a new scanner, task, age group, or intervention.

\subsection{Training and the transfer problem}
Repeated testing increases scores through item memory, strategy discovery, motivation, and familiarity. Adaptive training may produce genuine near-transfer while producing little far-transfer. The distinction matters because a task-specific learning curve can plateau or improve without changing a higher-order general factor. A strong experiment therefore uses novel stimuli, alternate forms, an active control, untrained transfer tasks, and delayed follow-up.

\subsection{Why the upper tail is especially difficult}
Standard batteries may not contain enough high-difficulty items to distinguish very high performers. A flat curve can result from a scoring ceiling, an imposed time limit, item exposure, or a true performance boundary. High-range item-response calibration is a prerequisite for interpreting a plateau. The calibration must also test differential item functioning across age, language, site, and education, because a numerically equal score may represent different latent performance.

\begin{table}[t]
\caption{Levels of the intelligence-limit question.}
\label{tab:levels}
\centering
\footnotesize
\setlength{\tabcolsep}{2pt}
\begin{tabular}{p{0.25\columnwidth}p{0.64\columnwidth}}
\toprule
Level & Scientific object and limit interpretation\\
\midrule
Score & Normed performance on one battery; may show a test or cohort artifact.\\
Domain ability & Latent reasoning, knowledge, memory, speed, or spatial factor; may have different trajectories.\\
General factor & Higher-order covariance among domains; requires invariance and model comparison.\\
Learning capacity & Improvement on new tasks; must separate practice, near transfer, and far transfer.\\
Biological constraint & Neural or genetic moderator of a trajectory; not automatically immutable.\\
\bottomrule
\end{tabular}
\end{table}

\section{Scientific Reframe and Hypotheses}
\label{sec:idea}
We call the proposed direction the \emph{CONSTRAINED-GENERAL-INTELLIGENCE-PLASTICITY-BENCHMARK}. The word limit is operationalized as a reproducible asymptote in a declared latent trajectory. The asymptote is conditional on task family, time constraint, age, health, environment, and available practice. A vector of domain-specific boundaries is the primary object; a general factor is a summary, not an assumed universal ruler.

\subsection{Hypotheses}
\textbf{H0: Measurement-equivalence hypothesis.} After high-range calibration and practice adjustment, some apparent global plateaus will relocate to particular domains or weaken. A universal plateau is not expected unless measurement invariance holds.

\textbf{H1: Environmental-shift hypothesis.} Education, sleep, health, and enrichment will shift latent performance heterogeneously, with larger effects in acquired or strategy-dependent domains. The shift will vary by age and baseline opportunity.

\textbf{H2: Conditional-asymptote hypothesis.} Within a fixed task family and resource condition, repeated learning curves will approach individual-specific asymptotes. These asymptotes will differ by domain, age, baseline ability, and health.

\textbf{H3: Transfer-limitation hypothesis.} Adaptive training will produce larger gains on near-transfer tasks than on far-transfer tasks. Evidence for a limit to general intelligence requires a plateau in cross-domain latent performance rather than a plateau in the trained task.

\textbf{H4: Network-mediation hypothesis.} Brain measures associated with processing efficiency and flexible network recruitment will partially mediate or moderate individual differences in the frontier. They will not fully determine the frontier, and compensation may preserve performance in some strata.

\subsection{Falsifiable outcomes}
If a plateau disappears after difficult items and time-unlimited forms are added, it is classified as a measurement ceiling. If a gain appears only on trained items or the short retest, it is classified as practice. If a latent cross-domain effect persists in the active-training arm at delayed follow-up, the result supports transferable plasticity. If a domain-specific slope approaches zero after calibration and remains stable across forms and sites, a conditional limit is supported for that domain. Heterogeneous asymptotes reject the claim that one IQ value describes all intelligence.

\section{Experiment Design}
\label{sec:design}
This protocol is design-only and requires ethics approval before implementation. It combines a longitudinal cohort, a randomized training module, and a neural sub-study. The planning target is approximately 1,200 participants across multiple sites and four age strata: adolescence, early adulthood, midlife, and later adulthood. Final enrollment should be set by preregistered power simulation, expected attrition, and the number of language and education strata. An intervention subset of approximately 400 participants is a planning value rather than a claimed sample.

\subsection{Sampling and controls}
Recruitment is stratified by education, sex, socioeconomic conditions, language background, and health status. Testing language, sensory accommodation, medication, chronic disease, sleep, and prior exposure to digital tasks are recorded. Age is not treated as a substitute for development, cohort, or health. The neurocognitive sub-study enrolls a consented subset and records scanner, motion, site, and preprocessing batch.

The measurement controls are alternate forms, high-range items, speeded and time-unlimited versions where valid, and a short retest subset. The intervention controls are an adaptive reasoning and working-memory curriculum versus a time-matched engaging non-adaptive active control. Both arms receive equivalent contact, feedback, and assessment schedules. Far-transfer tasks are never used as training material.

\subsection{Cognitive battery}
The battery contains fluid reasoning, crystallized knowledge, working memory, processing speed, spatial reasoning, and novel-task learning. Each domain contains items extending beyond ordinary screening ranges. A pilot calibrates item difficulty, discrimination, response time, and differential item functioning. Items with severe non-equivalence are revised or modeled before confirmatory analysis.

Speeded tasks are not interpreted as pure reasoning. Time-limit indicators and a speed factor are included in the measurement model. A person can gain accuracy while losing speed, or vice versa; the composite should not hide this tradeoff. Alternate forms prevent the same answer key from becoming a memory test.

\subsection{Environmental measurements}
Years and quality of education, sleep regularity, nutrition indicators that can be collected ethically, chronic disease burden, stress exposure, language, digital-task exposure, and opportunity measures are recorded longitudinally. These variables are moderators or explanatory covariates unless a design supports causal inference. The randomized training module supplies the principal causal test of a modifiable performance component.

\subsection{Neural measurements}
EEG estimates task-evoked timing, spectral dynamics, and trial-to-trial variability. Structural and resting-state MRI estimate candidate regional and network features. The analysis uses a reliability screen, predefined feature families, and nested cross-validation. Mediation is tested only with temporal ordering and confound sensitivity analyses. A neural predictor is not allowed to become a biological-destiny claim merely because it predicts a score in one sample.

\section{Formal Model and Analysis}
\label{sec:model}
For participant $i$, domain $d$, task or item $j$, and time $t$, observed performance is modeled as
\begin{equation}
y_{ijdt}=\nu_j+\lambda_j g_{it}+\delta_{jd}d_{idt}+\beta_j s_{it}+\rho_j p_{ijt}+\epsilon_{ijdt},
\label{eq:factor}
\end{equation}
where $g_{it}$ is a higher-order general factor, $d_{idt}$ is a domain factor, $s_{it}$ is a speed or time-limit component, and $p_{ijt}$ captures practice on the task family. Item parameters are estimated with an item-response model. Site, language, age, and batch terms are included hierarchically.

For a fixed task family, a conditional learning trajectory is represented by
\begin{equation}
g_i(t)=\theta_i-A_i\exp(-k_i t)+u_{i,t},
\label{eq:asymptote}
\end{equation}
where $\theta_i$ is an individual conditional asymptote, $A_i$ is available improvement, $k_i$ is a learning rate, and $u_{i,t}$ is time-varying deviation. The parameter $\theta_i$ is not a universal maximum. It is conditional on the item pool, time, health, environment, and training history.

A plateau is accepted only when four conditions are met. First, high-range item calibration shows that the scale still discriminates near the observed frontier. Second, the derivative of the trajectory remains within a preregistered small-slope interval over the declared horizon. Third, alternate forms and delayed follow-up reproduce the result. Fourth, the result is directionally stable at an independent site or cohort. Model fit alone is not sufficient because a flexible curve can mimic a plateau.

Measurement invariance is tested across age, site, language, and time. If scalar invariance fails, raw score comparisons and a universal IQ trend are not interpreted. The training effect is estimated by intent-to-treat mixed-effects analysis, with per-protocol analysis labeled sensitivity analysis. Technical repetitions estimate measurement error and do not increase the participant-level sample size.

For a neural sub-study, let $N_i$ be a vector of preregistered neural features and $Z_i$ denote environmental covariates. The prediction model is
\begin{equation}
\theta_i=\alpha+\eta^{\mathsf T}N_i+\zeta^{\mathsf T}Z_i+\omega_i,
\label{eq:neural}
\end{equation}
with nested cross-validation and site-held-out validation. A mediation analysis asks whether a neural feature transmits part of an intervention or environment association, but sensitivity analysis must quantify the effect of unmeasured confounding. Neither prediction nor mediation proves a hard ceiling.

\section{Construct Validation and Scale Extension}
\label{sec:validation}
The most important prerequisite for an upper-limit study is a scale that remains informative near the observed upper tail. A standard score can saturate before the person does. The pilot therefore fits an item-response model with a range of difficulty values that extends above the expected sample mean. For a binary item, a two-parameter form is
\begin{equation}
\Pr(y_{ij}=1\mid\vartheta_i)=[1+\exp\{-a_j(\vartheta_i-b_j)\}]^{-1},
\label{eq:irt}
\end{equation}
where $a_j$ is item discrimination, $b_j$ is item difficulty, and $\vartheta_i$ is the domain ability. The model is extended to polytomous or response-time outcomes when a binary score would discard useful information.

High-range calibration has three objectives. It must separate people who all answer ordinary items correctly, it must provide stable information at the relevant ability values, and it must keep the construct interpretable across forms. Items are not selected only because they are hard. An item that is difficult because of obscure cultural knowledge or language may reduce construct validity. The calibration records response time, strategy reports where appropriate, and differential item functioning so that a high score is not simply a measure of familiarity with one test culture.

Measurement invariance is evaluated in stages. Configural invariance asks whether the same domain pattern is present. Metric invariance asks whether factor loadings are comparable. Scalar invariance asks whether intercepts or thresholds permit latent mean comparisons. Strict invariance concerns residual scales and is desirable for some comparisons but not required for every latent trajectory. When full invariance fails, partial invariance or item-level modeling is declared before interpreting group or cohort differences. A failure of invariance is itself an important result: it limits the meaning of a numerical score and prevents an apparent age or cohort ceiling from being overgeneralized.

Construct validation also compares competing structures. A single common factor, a hierarchical general-factor model, a correlated-domain model, and a bifactor model are considered when justified by the battery. The preferred model is chosen using predictive performance, residual diagnostics, reliability, and interpretability rather than a single fit index. If the general factor is unstable across forms or ages, the study reports domain trajectories without forcing a general score. This is a stronger response to the limit question than silently treating model instability as measurement noise.

Speed and accuracy require a separate interpretation. A time limit may create a frontier that is primarily a motor or processing-speed frontier. Each speeded task is paired with a time-unlimited condition when the domain permits it, and accuracy, response time, and omission are modeled jointly. A participant who solves more difficult items slowly may have moved the reasoning frontier while not moving the speed frontier. Conversely, a faster response with no accuracy gain may reflect strategy or arousal. The analysis reports the tradeoff rather than collapsing it into one IQ number.

Novel-task learning is included because intelligence is partly expressed as adaptation to unfamiliar structure. At baseline, participants encounter a task family with no training history. The learning slope, strategy changes, and transfer to a structurally related but novel task are recorded. The asymptote in this module is conditional on the number of trials, feedback schedule, and task rule. It is not a direct estimate of all future learning. The module nevertheless distinguishes a person who begins high from a person who learns rapidly, which a single cross-sectional score cannot do.

The final construct-validation report includes item information curves, cross-form correlations, factor-loading stability, invariance results, high-tail precision, speed-accuracy tradeoffs, and the relationship between observed score and latent estimate. It also states where the scale is uninformative. A limit claim is suspended whenever the measurement information falls before the fitted asymptote, because the data then cannot distinguish a biological plateau from a scale ceiling.

\section{Analysis Deliverables and Reproducibility Path}
\label{sec:deliverables}
The project is organized as a sequence of locked deliverables. The first is a survey and source registry that records which claims concern scores, latent factors, brain measures, genes, or interventions. The second is a preregistered item-calibration and measurement plan. The third is the longitudinal analysis specification, including the age-cohort-environment decomposition, missing-data model, and plateau tolerance. The fourth is the randomized-training analysis plan with near- and far-transfer endpoints. The fifth is the neural analysis plan with reliability filtering and site-held-out validation.

Each deliverable has a stop rule. If high-range items are not reliable, no upper-tail conclusion is issued. If alternate forms are not equivalent, retest-adjusted trajectories are reported only as exploratory. If attrition is strongly performance-dependent and sensitivity intervals are too wide, the study reports non-identifiability rather than a stable limit. If the active control is not credible, training results are labeled expectancy-sensitive. If the neural sub-study cannot reach reliability or independent validation, it remains descriptive.

The final data package separates identifiers, raw responses, item parameters, factor scores, intervention assignments, environmental records, and imaging derivatives. Analysis scripts reproduce the primary tables from a versioned manifest. A public summary reports estimates and uncertainty without exposing personal scores or stigmatizing labels. This structure makes a null result valuable: future studies can determine whether the failed link was scale construction, cohort access, training transfer, or neural reliability.

The model classes and the failure ledger are reported together. The reader can therefore see whether a conclusion depends on the chosen curve, a particular site, a scoring transformation, or a restricted task family. A reproducible upper boundary is a result only when these dependencies are measured and bounded.

\section{Model Classes and Non-Identifiability}
\label{sec:modelclasses}
The answer depends on which of three model classes survives the preregistered checks. In a no-plateau model, the expected latent trajectory continues to change over the observed horizon, although its rate may become small. In a domain-specific plateau model, one or more domains have a stable asymptote while other domains continue to change. In a common-plateau model, a higher-order factor and most domains approach a shared boundary after measurement and practice are controlled. These classes have different implications and should not be collapsed into one p-value.

The comparison uses predictive performance, residual structure, and the coverage of the asymptote interval. A model that fits the mean but predicts poorly for the upper tail is inadequate for the limit question. A flexible spline can reproduce a flat segment without implying a true asymptote, while a rigid exponential can impose one by construction. The primary analysis therefore compares a bounded curve with a flexible non-bounded curve and checks whether the conclusion is stable when the functional form changes. Model averaging can be used when no single curve is clearly preferred, with the resulting uncertainty reported rather than hidden.

A conditional boundary is identifiable only over a declared support. If the highest-difficulty items remain easy for the top participants, the upper tail is not measured. If only one age cohort is observed, age and cohort are confounded. If all participants receive the same training and no active control is present, practice and intervention effects are confounded. If all data come from one language or one high-income site, the population scope is narrow. If neural scans are available only from the highest-performing volunteers, a brain-performance association can be selection-induced. These are conditions under which a universal limit cannot be estimated.

The analysis reports an explicit non-identifiability outcome. It is used when intervals for the trajectory derivative remain wide, when invariance fails for a large fraction of items, when attrition makes the missing-data sensitivity range too broad, or when the intervention control is not credible. Non-identifiability means that the data and design do not distinguish competing explanations. It is preferable to declaring either unlimited intelligence or a hard ceiling from an underpowered or saturated scale.

The frontier can also be represented as a distribution rather than a point. Let $\Theta_d$ denote the distribution of conditional asymptotes in domain $d$. The design estimates its mean, spread, and moderation by age, health, education, and baseline opportunity. A narrow $\Theta_d$ after extensive calibration would make a domain boundary more reproducible. A broad distribution would show that individual variation is the relevant object. A changing $\Theta_d$ across cohorts or interventions would show a movable population frontier. None of these outcomes implies that the distribution is fixed for all future environments.

This model-class framework keeps the historical claim about a thirty-point rise in perspective. Even if the observed normed score increase is replicated, it can coexist with a domain-specific or moving frontier. Conversely, even if a modern cohort shows little further score increase, a new item pool or time-unlimited task can reveal additional capacity. The inference must follow construct validation, not the rhetoric of either limitless growth or inevitable saturation.

\section{Separating Age, Cohort, and Environment}
\label{sec:decomposition}
Cross-sectional age differences are not the same as developmental change. People born in different years have different schools, nutrition, technology, disease exposures, and selective-survival histories. A lower score in an older group could reflect aging, a cohort difference, a health difference, or a measurement problem. The design therefore combines age strata with repeated follow-up and measures the environments that shape opportunities.

Let $a$ denote age, $c$ denote birth cohort, and $e$ denote a vector of environmental exposures. A decomposition model for a domain factor is
\begin{equation}
d_{idt}=\mu_d+f_d(a_{it})+h_d(c_i)+q_d^{\mathsf T}e_{it}+b_{site(i)}+r_{it},
\label{eq:ace}
\end{equation}
where $f_d$ is an age trajectory, $h_d$ is a cohort term, and $q_d$ describes measured environmental moderation. The model is not used to claim that age, cohort, and environment can always be perfectly separated; their partial collinearity is an empirical limitation. Its purpose is to prevent an age curve from being reported as a universal biological law.

Longitudinal observations identify within-person change better than a single cross-section, but they also introduce retest effects and selective attrition. Alternate forms, a short retest calibration, and explicit practice terms are essential. The cohort component is estimated from overlapping birth cohorts and sites rather than from a single generational contrast. If a common age curve fits poorly while domain-specific curves are stable, the conclusion is a heterogeneous frontier.

Environmental measurements are interpreted with care. Education quality is not equivalent to years of schooling, and sleep, nutrition, and health are partly consequences of social conditions. A covariate association can therefore be informative without being causal. The randomized training module supplies a sharper causal contrast for one modifiable input, while the broader environmental variables define the conditions under which a conditional frontier is expressed. Interaction terms test whether the training response differs by baseline ability, age, health, or educational opportunity.

The same decomposition applies to the historical Flynn effect. A cohort gain in a crystallized knowledge factor may differ from a gain in fluid reasoning or processing speed. A score trend that disappears after item calibration is a measurement result. A latent trend that remains after invariance testing is evidence of a population-level change in the measured construct, but it is still not evidence of an unbounded capacity. The appropriate statement is always indexed by domain and condition.

\section{Power, Uncertainty, and Missingness}
\label{sec:power}
The planned sample size is not treated as evidence of precision. Before data collection, simulations should vary factor loadings, item discrimination, attrition, site heterogeneity, training effects, and the width of the high-range item pool. The simulations should quantify false detection of a plateau, false rejection of a plateau, coverage of the asymptote interval, and power to detect a delayed cross-domain training effect. A sample is adequate only for the estimands that its simulated precision supports.

The primary plateau decision is an interval decision rather than a test of whether a slope is exactly zero. Let $m_d(t)$ be the expected slope of domain $d$. A practical plateau criterion is that the simultaneous interval for $m_d(t)$ lies within $[-\delta_d,\delta_d]$ for a prespecified horizon, after item calibration and practice adjustment. The tolerance $\delta_d$ has to be expressed in interpretable latent units and justified before seeing the data. A very small tolerance can make every realistic study inconclusive; a large tolerance can declare a limit too easily.

Uncertainty is propagated from item parameters to factor scores and from factor scores to individual trajectories. Bootstrap or posterior predictive checks assess whether the fitted asymptote is driven by a few high performers or by a particular site. The report includes calibration uncertainty, parameter correlation, sensitivity to priors or penalization, and predictive performance in held-out forms. A high point estimate with a wide interval is not a reliable upper boundary.

Missingness is modeled at the measurement occasion and participant levels. Reasons for missing data include illness, fatigue, access, scanner motion, language difficulty, technical failure, and withdrawal. The primary analysis uses a declared missing-data model, with pattern-mixture and inverse-probability sensitivity analyses. Complete-case estimates are not treated as primary because dropout may depend on performance or intervention burden. The study also reports whether the highest-performing participants are disproportionately missing difficult items because of time limits or boredom.

Multiple comparisons are controlled by hierarchical ordering. Confirmatory claims concern the general factor, domain asymptotes, measurement invariance, and delayed intervention transfer. Neural features, subgroups, and environmental interactions are secondary or exploratory unless preregistered. Exploratory discoveries can motivate new studies but cannot overrule a failed primary gate. This keeps the answer to a limit question from being selected from dozens of plausible curves after the fact.

\section{Robustness Checks and Failure Ledger}
\label{sec:robustness}
The central failure mode is circularity: defining an intelligence scale, observing a plateau on that scale, and then calling the scale's endpoint a human limit. The benchmark counters this by testing more difficult items, alternate forms, time-unlimited versions, and tasks from different domains. A plateau that survives only one scoring rule is not accepted.

The second failure mode is confusing a higher mean with a higher maximum. Environmental improvement may move the lower and middle parts of a distribution while leaving the upper tail unchanged. Conversely, a small average shift may conceal a large shift in a disadvantaged subgroup. Distributional analyses therefore report quantiles, variance, rank changes, and ceiling sensitivity rather than only the mean factor score. The primary limit claim remains conditional on the observed support of the sample.

The third failure mode is treating a neural association as a mechanism. A larger or more efficient network measure may be a consequence of learning, a correlate of socioeconomic opportunity, or a compensatory response to age-related change. The neural sub-study reports temporal ordering, cross-validation, site-held-out performance, and confounding sensitivity. It does not use a brain measure to assign a fixed score to an individual or to infer that an intervention cannot work.

The fourth failure mode is unmeasured tool use. Humans solve problems with language, writing, calculators, computers, collaborators, and cultural knowledge. The design records tool availability and defines whether a task is tool-free, tool-assisted, or socially supported. A limit on unaided matrix reasoning is not a limit on problem solving in a technologically extended environment. This distinction is essential when the question is translated from a test room to real-world intelligence.

Each stage maintains a failure ledger. The ledger records items removed for differential functioning, forms that fail equivalence, sites with unreliable timing, training sessions below adherence, scanner runs with excessive motion, and endpoints that fail safety or reliability criteria. The final report distinguishes not tested, not identifiable, null under the declared design, and negative evidence against the hypothesis. This vocabulary prevents a missing measurement from being silently interpreted as a zero effect.

Replication is staged. First replicate the psychometric model within an independent item pool. Then replicate the latent trajectory at an independent site or cohort. Finally test whether the intervention effect is stable under a different trainer, language, or delivery format. If the asymptote changes across these checks, the correct result is a context-dependent boundary, not a failed study or a universal ceiling.

\section{Staged Protocol and Decision Gates}
\label{sec:gates}
\subsection{S0: calibration and pilot}
Verify high-range coverage, alternate-form equivalence, test-retest reliability, response-time instrumentation, and neural preprocessing reliability. The pilot must include participants near the upper observed range. A scoring rule is frozen before confirmatory data are inspected.

\subsection{S1: baseline cross-sectional measurement}
Estimate domain factors and the higher-order factor across age and environmental strata. Compare no-plateau, domain-specific-plateau, and common-plateau models using predictive calibration, posterior or confidence intervals, residual checks, and invariance diagnostics. A common curve is not preferred solely because it is simpler.

\subsection{S2: longitudinal follow-up}
Repeat alternate forms at preregistered intervals. Estimate individual trajectories, retest effects, attrition mechanisms, cohort-by-age interactions, and time-varying health. The missing-data plan distinguishes dropout related to fatigue, illness, access, and performance.

\subsection{S3: randomized plasticity module}
Compare adaptive training with the active control at post-test and delayed follow-up. The primary contrast is change in latent domain factors and the higher-order factor. Near-transfer shares component processes with training but uses new stimuli. Far-transfer is structurally distinct. A trained-task gain without a latent cross-domain gain is not evidence that general intelligence increased.

\subsection{S4: neural mediation and external validation}
Test only the predefined neural feature families and validate the measurement model at an independent site or cohort. If a neural signature predicts the baseline asymptote but fails to predict intervention response, it is a correlate or constraint candidate, not a complete mechanism. If the signature fails under a new scanner or language, its scope is narrowed.

\begin{table}[t]
\caption{Decision gates for interpreting an apparent intelligence limit.}
\label{tab:gates}
\centering
\footnotesize
\setlength{\tabcolsep}{2pt}
\begin{tabular}{p{0.27\columnwidth}p{0.62\columnwidth}}
\toprule
Evidence pattern & Permitted interpretation\\
\midrule
Plateau removed by hard items & Psychometric ceiling; no biological limit inferred.\\
Gain only on trained items & Practice or near-task learning; no general transfer.\\
Gain after active control on untrained domains & Transferable plasticity is supported.\\
Stable domain slope near zero & Conditional domain asymptote is supported.\\
Stable asymptote across sites and forms & Reproducibility of the conditional boundary is supported.\\
Different domains show different asymptotes & Intelligence is better represented as a frontier vector.\\
Neural correlate predicts but does not mediate change & Correlate or moderator; not a fixed destiny.\\
\bottomrule
\end{tabular}
\end{table}

\section{Threats to Validity, Ethics, and Reproducibility}
\label{sec:risks}
The largest threat is construct drift. If item content changes, an apparent cohort gain or decline may reflect the tested construct rather than a population change. Measurement invariance and item-response calibration are therefore primary, not optional. A second threat is selective participation. Highly educated, healthy, or technologically comfortable volunteers can make an upper-tail estimate appear more stable than it is. Stratified recruitment and transparent missingness analyses are needed.

Attrition can bias learning curves, especially when fatigue, illness, or poor access is related to performance. The protocol retains reasons for dropout and reports sensitivity analyses. Site, language, scanner, and operator effects are modeled as hierarchical sources of variation. An independent validation cohort is more informative than a second analysis of the same participants.

The intervention raises a different threat: expectancy. The active control is matched for contact and engagement, and the training hypothesis is stated before outcomes are inspected. A short-term improvement can be caused by strategy adoption, motivation, or item familiarity. Delayed follow-up and far-transfer tasks are necessary, but even they do not show that all aspects of intelligence changed.

Participant welfare and privacy are central. Consent or age-appropriate assent, withdrawal without penalty, secure separation of identifiers, and MRI/EEG safety screening are required. Results should not be described as fixed personal worth or immutable ability. Genetic data are not required for the core study; a future genetic extension would require separate consent, governance, and careful limits on individual prediction. Cognitive training is low risk but should include fatigue monitoring and stopping rules.

Reproducibility requires release of the preregistration, item-calibration report, form-construction rules, analysis code, model-checking plots, and de-identified summaries where permitted. Failed calibration runs, excluded items, scanner failures, and negative transfer results remain in an audit trail. Technical replication does not substitute for participant-level replication.

\section{Interpretation, Limitations, and Forecast}
\label{sec:forecast}
The design yields several scientifically distinct conclusions. A stable latent asymptote after high-range calibration supports a conditional limit for a specified ability. A training shift in latent cross-domain performance supports plasticity of that frontier. A domain-specific pattern rejects the idea that one IQ value is a universal maximum. An environmental moderation effect shows that the frontier is expressed through opportunity and state, not only biology. A neural association that fails to predict change is not a complete explanation.

The design cannot prove that no future educational technology, social environment, or tool-assisted system could raise performance. It cannot identify an absolute maximum for every task or person. The proposed cohort may still underrepresent rare cognitive profiles, and any MRI/EEG sub-study will be smaller than the psychometric cohort. These limitations bound the claim rather than invalidate the benchmark.

The most defensible forecast is conditional. Human cognition likely has domain- and condition-specific constraints arising from information processing, memory, speed, development, health, and resource allocation. Those constraints need not be fixed at one value: education, health, sleep, and learning can move performance and alter the rate at which a person approaches a task frontier. A universal IQ ceiling is not supported by cohort gains, because those gains are heterogeneous and norm-dependent. Nor is unlimited general intelligence established by the existence of new problems, because a new problem can be solved with tools, accumulated culture, collaboration, or domain-specific expertise.

The practical answer is therefore not a binary yes-or-no answer in the abstract. It is: specify the ability, build a scale that does not saturate, separate practice from transfer, estimate the trajectory, and test whether the asymptote moves under a controlled intervention. Only then can a limit be discussed as a scientific property rather than a slogan.

\section{Conclusion}
\label{sec:conclusion}
IQ scores can rise across generations, and structured experience can improve measured performance, without implying that a single biological quantity increases without bound. The proposed conditional intelligence frontier benchmark treats a limit as a reproducible asymptote in a calibrated, domain-specific, and resource-defined latent trajectory. It combines high-range items, alternate forms, longitudinal measurement, active-control training, far-transfer tasks, and a carefully bounded neural sub-study.

A convincing limit claim requires measurement invariance, practice separation, delayed replication, and cross-site stability. A convincing plasticity claim requires change beyond an active control on untrained tasks. A neural correlate may explain variation or compensation but is not an immutable destiny. This report presents a design-only protocol and no participant or intervention result. Its forecast is that human intelligence is neither one unrestricted scalar nor one already-measured universal ceiling; it is a heterogeneous frontier that can be constrained, expressed, and partly shifted under specific conditions.

\begin{thebibliography}{00}
\bibitem{pietschnig2015} J. Pietschnig and M. Voracek, One century of global IQ gains: A formal meta-analysis of the Flynn effect (1909--2013), \emph{Perspectives on Psychological Science}, vol. 10, no. 3, pp. 282--306, 2015, doi: 10.1177/1745691615577701.
\bibitem{ritchie2018} S. J. Ritchie and E. M. Tucker-Drob, How much does education improve intelligence? A meta-analysis, \emph{Psychological Science}, vol. 29, no. 8, pp. 1358--1369, 2018, doi: 10.1177/0956797618774253.
\bibitem{plomin2018} R. Plomin and S. von Stumm, The new genetics of intelligence, \emph{Nature Reviews Genetics}, vol. 19, pp. 148--159, 2018, doi: 10.1038/nrg.2017.105.
\bibitem{savage2018} J. E. Savage \emph{et al.}, Genome-wide association study of intelligence: A meta-analysis of 269,867 individuals and 1,016,443 SNPs, \emph{Molecular Psychiatry}, vol. 23, pp. 1169--1181, 2018, doi: 10.1038/s41588-018-0147-3.
\bibitem{deary2010} I. J. Deary, L. Penke, and W. Johnson, The neuroscience of human intelligence differences, \emph{Nature Reviews Neuroscience}, vol. 11, pp. 201--211, 2010, doi: 10.1038/nrn2793.
\bibitem{deary2021} I. J. Deary \emph{et al.}, Genetic variation, brain, and intelligence differences, \emph{Molecular Psychiatry}, vol. 26, pp. 4047--4061, 2021, doi: 10.1038/s41380-021-01027-y.
\bibitem{nisbett2012} R. E. Nisbett \emph{et al.}, Intelligence: New findings and theoretical developments, \emph{American Psychologist}, vol. 67, no. 2, pp. 130--159, 2012, doi: 10.1037/a0026699.
\bibitem{carroll1993} J. B. Carroll, \emph{Human Cognitive Abilities: A Survey of Factor-Analytic Studies}. Cambridge, U.K.: Cambridge University Press, 1993.
\end{thebibliography}

\end{document}
